\hypertarget{chapter-06-page-129}{}
\subsection{BMO 空间}
\label{sec:ch6-bmo}

给定函数 $f\in L^1_{\mathrm{loc}}(\mathbb R^n)$ 和立方体 $Q$, 用 $f_Q$ 表示 $f$ 在 $Q$ 上的平均值:
\[
 f_Q=\frac1{|Q|}\int_Qf.
\]
定义锐极大函数
\begin{equation}
\label{eq:ch6-sharp-maximal-function}
 M^\#f(x)=\sup_{Q\ni x}\frac1{|Q|}\int_Q|f-f_Q|,
\end{equation}
其中上确界取遍所有包含 $x$ 的立方体 $Q$. 每个积分都度量 $f$ 在立方体 $Q$ 上的平均振荡. 若函数 $M^\#f$ 有界, 就说 $f$ 具有有界平均振荡. 具有此性质的函数空间记作 BMO:
\[
 \mathrm{BMO}=\{f\in L^1_{\mathrm{loc}}:M^\#f\in L^\infty\}.
\]
在 BMO 上定义
\[
 \|f\|_*=\|M^\#f\|_\infty.
\]
严格说来这不是范数, 因为几乎处处为常数的函数具有零振荡. 不过, 只有这些函数具有该性质, 所以通常把 BMO 看作上述空间关于常数函数空间的商空间. 换言之, 相差一个常数的两个函数在 BMO 中视为同一函数. 在此空间上, $\|\cdot\|_*$ 是范数, 且 BMO 是 Banach 空间.

容易看出逐点不等式
\[
 M^\#f(x)\le C_nMf(x),
\]
其中 $M$ 是 Hardy--Littlewood 极大函数. 常数依赖维数 $n$, 但若用 $M''$ 代替 $M$, 即取遍所有包含 $x$ 的立方体的极大算子, 则常数为 $2$.

\begin{Proposition}
\label{prop:ch6-bmo-equivalent-oscillation}
有
\begin{equation}
\label{eq:ch6-bmo-infimum-equivalence}
 \frac12\|f\|_*
 \le\sup_Q\inf_{a\in\mathbb C}\frac1{|Q|}\int_Q|f(x)-a|\,dx
 \le\|f\|_*,
\end{equation}
以及
\begin{equation}
\label{eq:ch6-absolute-value-sharp-maximal}
 M^\#(|f|)(x)\le2M^\#f(x).
\end{equation}
\end{Proposition}

\hypertarget{chapter-06-page-130}{}
\begin{Proof}
\eqref{eq:ch6-bmo-infimum-equivalence} 中第二个不等式是直接的. 为证明第一个不等式, 注意对每个 $a$,
\[
 \int_Q|f(x)-f_Q|\,dx
 \le\int_Q|f(x)-a|\,dx+\int_Q|a-f_Q|\,dx
 \le2\int_Q|f(x)-a|\,dx.
\]
两边除以 $|Q|$, 先对 $a$ 取下确界, 再对 $Q$ 取上确界即可.

在 \eqref{eq:ch6-bmo-infimum-equivalence} 中取 $a=|f_Q|$, 并利用
\[
 \frac1{|Q|}\int_Q\bigl||f(x)|-|f_Q|\bigr|\,dx
 \le\frac1{|Q|}\int_Q|f(x)-f_Q|\,dx,
\]
即得 \eqref{eq:ch6-absolute-value-sharp-maximal}.
\end{Proof}

\eqref{eq:ch6-bmo-infimum-equivalence} 定义了一个与 $\|\cdot\|_*$ 等价的范数, 它给出一种无需使用 $f$ 在 $Q$ 上的平均值便可证明 $f\in\mathrm{BMO}$ 的方法: 只需找到一个可依赖于 $Q$ 的常数 $a$, 使
\[
 \frac1{|Q|}\int_Q|f(x)-a|\,dx\le C,
\]
其中 $C$ 与 $Q$ 无关.

由 \eqref{eq:ch6-absolute-value-sharp-maximal}, 若 $f\in\mathrm{BMO}$, 则 $|f|\in\mathrm{BMO}$. 反命题不成立. 这印证了定义所暗示的事实: 属于 BMO 不只是大小问题. 显然 $L^\infty\subset\mathrm{BMO}$, 但 BMO 中也有无界函数. $\mathbb R$ 上的典型例子是
\[
 f(x)=
 \begin{cases}
  \log(1/|x|),&|x|<1,\\
  0,&|x|\ge1.
 \end{cases}
\]
但容易看出 $\operatorname{sgn}(x)f(x)\notin\mathrm{BMO}$, 尽管 $f$ 是它的绝对值.

下面的结果说明 BMO 与奇异积分之间的联系.

\begin{Theorem}
\label{thm:ch6-linfty-to-bmo}
设 $T$ 满足定理~\ref{thm:ch5-generalized-calderon-zygmund-bounds} 的假设, 且其核满足条件 \eqref{eq:ch5-generalized-hormander-x}. 若 $f$ 是具有紧支集的有界函数, 则 $Tf\in\mathrm{BMO}$, 且
\[
 \|Tf\|_*\le C\|f\|_\infty.
\]
\end{Theorem}

\begin{Proof}
固定中心为 $c_Q$ 的立方体 $Q\subset\mathbb R^n$, 并令 $Q^*$ 是中心为 $c_Q$, 边长为 $Q$ 的 $2\sqrt n$ 倍的立方体. 将 $f$ 分解为 $f=f_1+f_2$, 其中 $f_1=f\chi_{Q^*}$.

\hypertarget{chapter-06-page-131}{}
令 $a=Tf_2(c_Q)$. 由条件 \eqref{eq:ch5-generalized-hormander-x} 以及 $T$ 在 $L^2$ 上有界,
\[
\begin{aligned}
 \frac1{|Q|}\int_Q|Tf(x)-a|\,dx
 &\le\frac1{|Q|}\int_Q|Tf_1(x)|\,dx
 +\frac1{|Q|}\int_Q|Tf_2(x)-Tf_2(c_Q)|\,dx\\
 &\le\left(\frac1{|Q|}\int_Q|Tf_1(x)|^2\,dx\right)^{1/2}\\
 &\quad+\frac1{|Q|}\int_Q\left|\int_{\mathbb R^n\setminus Q^*}
 [K(x,y)-K(c_Q,y)]f(y)\,dy\right|dx\\
 &\le C\left(\frac1{|Q|}\int_{Q^*}|f(x)|^2\,dx\right)^{1/2}\\
 &\quad+\frac1{|Q|}\int_Q\int_{\mathbb R^n\setminus Q^*}
 |K(x,y)-K(c_Q,y)|\,dy\,dx\,\|f\|_\infty\\
 &\le C\|f\|_\infty.
\end{aligned}
\]
\end{Proof}

定理~\ref{thm:ch6-linfty-to-bmo} 是寻找一个包含奇异积分作用于有界函数所得像的空间的第一步. 但是, 具有紧支集的有界函数在 $L^\infty$ 中不稠密, 因而不能通过连续性把 $T$ 延拓到整个空间. 我们改为把 $T$ 的定义延拓到整个 $L^\infty$.

设 $f$ 是有界函数, $Q$ 是 $\mathbb R^n$ 中以原点为中心的立方体. 如上定义 $Q^*$, 再次写 $f=f_1+f_2$, 其中 $f_1=f\chi_{Q^*}$. 因为 $f_1$ 有界且具有紧支集, $Tf_1$ 作为 $L^2$ 函数有定义, 所以 $Tf_1(x)$ 对几乎每个 $x$ 存在. 对 $x\in Q$ 定义
\begin{equation}
\label{eq:ch6-linfty-extension}
 Tf(x)=Tf_1(x)+\int_{\mathbb R^n}[K(x,y)-K(0,y)]f_2(y)\,dy.
\end{equation}
若 $K$ 满足 \eqref{eq:ch5-generalized-hormander-x}, 则该积分收敛, 因为其绝对值不超过
\[
 \|f\|_\infty\int_{\mathbb R^n\setminus Q^*}|K(x,y)-K(0,y)|\,dy.
\]

现在设 $\widetilde Q$ 是另一个包含 $Q$ 且以原点为中心的立方体. 对 $x\in Q$, 我们得到 $Tf(x)$ 的两个定义. 令 $\widetilde f_1=f\chi_{\widetilde Q^*}$, $\widetilde f_2=f\chi_{\mathbb R^n\setminus\widetilde Q^*}$, 则两个定义之差为
\[
\begin{aligned}
 T(f_1-\widetilde f_1)(x)
 &+\int_{\mathbb R^n\setminus Q^*}[K(x,y)-K(0,y)]f(y)\,dy\\
 &-\int_{\mathbb R^n\setminus\widetilde Q^*}[K(x,y)-K(0,y)]f(y)\,dy
 =-\int_{\widetilde Q^*\setminus Q^*}K(0,y)f(y)\,dy.
\end{aligned}
\]

\hypertarget{chapter-06-page-132}{}
右端与 $x$ 无关. 因此, 两个定义作为 BMO 函数相同, 因为 BMO 中相差常数的函数视为同一函数. 于是可以用 \eqref{eq:ch6-linfty-extension} 定义 $Tf$. 进一步, 沿用定理~\ref{thm:ch6-linfty-to-bmo} 的论证可证, 若 $f\in L^\infty$, 则 $Tf\in\mathrm{BMO}$. 还应注意, 这里不必只取以原点为中心的立方体, 任意立方体都可以.

若 $f$ 是具有紧支集的有界函数, 则对充分大的 $Q$ 有 $f=f_1$, 所以 \eqref{eq:ch6-linfty-extension} 与原定义一致. 更一般地, 若 $f\in\mathcal S$, 则两个定义作为 BMO 函数一致: 因为
\[
 \int_{\mathbb R^n\setminus Q^*}K(0,y)f(y)\,dy
\]
绝对收敛到一个与 $x$ 无关的值, 两个定义只相差一个常数.

\begin{Example}
\label{ex:ch6-hilbert-transform-sign}
令 $f(x)=\operatorname{sgn}(x)$. 我们求它在 Hilbert 变换下的像, 其核为 $K(x,y)=[\pi(x-y)]^{-1}$. 设 $|x|<a/2$, 则
\[
\begin{aligned}
 \pi Hf(x)
 &=\operatorname{p.v.}\int_{-a}^a\frac{\operatorname{sgn}(y)}{x-y}\,dy
 +\lim_{N\to\infty}\left(\int_{-N}^{-a}+\int_a^N\right)
 \left(\frac1{x-y}+\frac1y\right)\operatorname{sgn}(y)\,dy\\
 &=2\log|x|-2\log(a).
\end{aligned}
\]
因此忽略常数后,
\[
 H(\operatorname{sgn}(x))=\frac2\pi\log|x|.
\]
这间接证明了 $\log|x|\in\mathrm{BMO}$. 该结果也与 Hilbert 变换最初的动机相符: 存在上半平面上的解析函数, 当趋近实轴时, 其实部趋于 $\operatorname{sgn}(x)$, 虚部趋于 $\frac2\pi\log|x|$. 这样的函数可以显式写成
\[
 i\log(iz)=i\log(-y+ix)=i\log(x^2+y^2)^{1/2}+\arctan(x/y).
\]
当 $y\to0$ 时, 该函数的极限为
\[
 \frac\pi2\operatorname{sgn}(x)+i\log|x|.
\]
\end{Example}
